\subsection{具有零微分的多项式的刻画}

命题 6. --- 设 K 为一个交换环，A 为多项式代数 \(K[X_{i}]_{i \in I}\)，S 为 A 中满足 dF = 0 的元素 F 的集合。

\begin{enumerate}
\def\labelenumi{\alph{enumi})}
\item
  如果 \(\mathbf{K}\) 是特征为 0 的环，则 \(\mathbf{S} = \mathbf{K}\) 。
\item
  如果 \(\mathbf{K}\) 是特征为 \(p \neq 0\) 的环，则 \(\mathbf{S} = \mathbf{K}[\mathbf{X}_i^p]_{i \in I}\)；此外如果 \(\mathbf{K}\) 是完全环，则 \(\mathbf{K} = \mathbf{A}^p\) 。
\end{enumerate}

映射 \(F \mapsto dF\)（从 A 到 A 的 K-微分形式模 \(\Omega_{\mathbf{K}}(\mathbf{A})\) 中）是 K-线性的，并且满足关系

\[
d \left(\mathrm{FF} ^ {\prime}\right) = \mathrm{F} \cdot d \mathrm{F} ^ {\prime} + \mathrm{F} ^ {\prime} \cdot d \mathrm{F}
\]

（III，第 569 页）。因此 S 是一个子代数。

当 K 的特征为 \(p \neq 0\) 时，令 \(T = K[X_{i}^{p}]_{i \in I}\)；于是若 K 是完全的，我们有 \(T = A^{p}\)（V，第 4 页，命题 2）；此外，我们有 \(d(X_{i}^{p}) = pX_{i}^{p-1}\)，\(dX_{i} = 0\) 对所有 \(i \in I\) 成立，因此 A 的子代数 S 包含 T。若 K 的特征为 0，我们令 T = K，并且仍然得到 \(T \subset S\) 。还需证明 S 包含于 T。

对每个有限子集 \(J\) （属于 \(I\)），令 \(A_{J}\) 为 \(A\) 的由族 \((X_{j})_{j\in J}\) 生成的子代数。我们有 \(A_{\emptyset} = K\) 且 \(A = \bigcup_{J\in I}A_{J}\)；因此只需证明

关系 \(S \cap A_{J} \subset T\)，我们将通过对 J 的基数进行归纳来证明。因此，设 J 是 I 的一个有限子集，使得 \(S \cap A_{J} \subset T\)，设 i 是 I - J 的一个元素，且 \(J' = J \cup \{i\}\) 。\(A_{J}\) 的每个元素 F 都可以以唯一的方式写成形式

\[
\mathrm{F} = \sum_ {n = 0} ^ {\infty} \mathrm{F} _ {n} \cdot \mathrm{X} _ {i} ^ {n},\tag{5}
\]

其中 \(\mathbf{F}_n\in \mathbf{A}_{\mathrm{J}}\) 对所有 \(n\geqslant 0\)，然后

\[
d \mathrm{F} = \sum_ {n = 0} ^ {\infty} \mathrm{X} _ {i} ^ {n} \cdot d \mathrm{F} _ {n} + \sum_ {n = 0} ^ {\infty} n \mathrm{X} _ {i} ^ {n - 1} \mathrm{F} _ {n} \cdot d \mathrm{X} _ {i}.\tag{6}
\]

假设 F 属于 S；族 \((d\mathbf{X}_{r})_{r\in I}\) 是 A-模 \(\Omega_{\mathrm{K}}(\mathbf{A})\) 的一组基（III，第 570 页），且 \(dF_{n}\) 是微分 \(dX_{j}\)（对 \(j\in J\)）的线性组合，因为 \(F_{n}\in A_{J}=K[X_{j}]_{j\in J}\) 。由 (6) 我们则有 \(dF_{n}=0\) 且 \(nF_{n}=0\)，对每个整数 \(n\geqslant0\)。由归纳假设，\(F_{n}\in T\) 对所有 n 成立，因为 \(dF_{n}=0\) 。

\begin{enumerate}
\def\labelenumi{\alph{enumi})}
\item
  若 K 的特征为 0，则我们有 \(nF_{n} = 0\) 对所有 \(n \geqslant 1\) 成立，从而由命题 1（V，第 2 页）得 \(F_{n} = 0\)；因此我们有 \(F = F_{0}\)，从而 \(F \in T\) 。
\item
  如果 \(\mathbf{K}\) 的特征为 \(p \neq 0\)，则 \(\mathbf{A}\) 是域 \(\mathbf{F}_p\) 上的代数，且关系 \(n\mathrm{F}_n = 0\) 蕴含 \(\mathbf{F}_n = 0\)，对每个整数 \(n\) （不被 \(p\) 整除）成立。因此我们有 \(\mathbf{F} = \sum_{m=0}^{\infty} \mathbf{F}_{mp} \mathbf{X}_i^{mp}\)，从而 \(\mathbf{F} \in \mathbf{T}\) 。
\end{enumerate}

注记。--- 当 K 的加法群无挠时，我们仍有 S = K；这可由上述证明或 V 第 3 页的注记 5 得出。

推论。--- 设 K 为一个域，F(X) 是系数在 K 中的多项式，其导数 F'(X) 为零。

\begin{enumerate}
\def\labelenumi{\alph{enumi})}
\item
  如果 \(\mathbf{K}\) 的特征为 0，则 \(\mathbf{F} \in \mathbf{K}\) 。
\item
  如果 \(\mathbf{K}\) 的特征为 \(p \neq 0\)，则存在一个多项式 \(\mathbf{G}(\mathbf{X})\)，使得 \(\mathbf{F}(\mathbf{X}) = \mathbf{G}(\mathbf{X}^p)\) 。
\end{enumerate}

因为我们有 \(d\mathbf{F} = \mathbf{F}' \cdot d\mathbf{X} = 0\) 。

